\documentclass[12pt]{article}

\usepackage[utf8]{inputenc}
\usepackage{amsmath, amssymb, amsthm}
\usepackage{geometry}
\geometry{a4paper, margin=2.5cm}

\title{Puncte Singulare Izolate ale Funcțiilor Olomorfe}
\author{}
\date{}

\theoremstyle{definition}
\newtheorem{definition}{Definiție}

\theoremstyle{plain}
\newtheorem{proposition}{Propoziție}

\begin{document}

\maketitle

\begin{definition}
Fie $\Omega \subset \mathbb{C}$ o mulțime deschisă și $z_0 \in \Omega$. Spunem că $z_0$ este
\emph{punct singular izolat} pentru o funcție $f : \Omega \setminus \{z_0\} \to \mathbb{C}$ dacă
există $r > 0$ astfel încât $f$ este olomorfă în domeniul perforat


\[
0 < |z - z_0| < r.
\]


\end{definition}

\begin{definition}
Fie $z_0$ un punct singular izolat al funcției $f$. Clasificarea singularității este:
\begin{enumerate}
    \item[(i)] \emph{Singularitate eliminabilă}: dacă există o funcție olomorfă $\tilde f$ definită
    în $|z - z_0| < r$ astfel încât $\tilde f(z) = f(z)$ pentru $z \ne z_0$.

    \item[(ii)] \emph{Pol de ordin $m$}: dacă există $m \ge 1$ astfel încât funcția
    

\[
    (z - z_0)^m f(z)
    \]


    se extinde la o funcție olomorfă nenulă în $z_0$.

    \item[(iii)] \emph{Singularitate esențială}: dacă singularitatea nu este nici eliminabilă,
    nici pol; echivalent, seria Laurent a lui $f$ în jurul lui $z_0$ conține o infinitate de termeni negativi.
\end{enumerate}
\end{definition}

\begin{proposition}
Fie $z_0$ un punct singular izolat al funcției $f$. Atunci:
\begin{enumerate}
    \item[(i)] $z_0$ este o singularitate eliminabilă dacă și numai dacă
    

\[
    \lim_{z \to z_0} f(z) \ \text{există și este finită}.
    \]



    \item[(ii)] $z_0$ este un pol dacă și numai dacă
    

\[
    \lim_{z \to z_0} |f(z)| = \infty.
    \]



    \item[(iii)] $z_0$ este o singularitate esențială dacă și numai dacă limitele de mai sus nu există,
    iar $f$ nu este mărginită în nicio vecinătate perforată a lui $z_0$.
\end{enumerate}
\end{proposition}

\begin{definition}
Fie $f$ o funcție olomorfă în afara unui disc suficient de mare, adică există $R > 0$ astfel încât


\[
f \text{ este olomorfă în } \{ z \in \mathbb{C} : |z| > R \}.
\]


Spunem că \emph{infinitul} este un punct singular izolat al lui $f$ dacă funcția


\[
g(z) = f\!\left(\frac{1}{z}\right)
\]


are un punct singular izolat în $z = 0$.
\end{definition}

\begin{definition}
Clasificarea singularității lui $f$ la infinit se face prin funcția $g(z) = f(1/z)$:
\begin{enumerate}
    \item[(i)] \emph{Infinitul este o singularitate eliminabilă} dacă $g$ are o singularitate eliminabilă în $0$, adică
    

\[
    \lim_{z \to 0} g(z) \ \text{există și este finită}.
    \]



    \item[(ii)] \emph{Infinitul este un pol de ordin $m$} dacă $g$ are un pol de ordin $m$ în $0$, adică
    

\[
    \lim_{z \to 0} |g(z)| = \infty.
    \]



    \item[(iii)] \emph{Infinitul este o singularitate esențială} dacă $g$ are o singularitate esențială în $0$, adică
    limitele de mai sus nu există, iar $g$ nu este mărginită în nicio vecinătate perforată a lui $0$.
\end{enumerate}
\end{definition}

\begin{proposition}
Pentru o funcție olomorfă $f$ în afara unui disc suficient de mare, următoarele caracterizări sunt valabile:
\begin{enumerate}
    \item[(i)] Infinitul este o singularitate eliminabilă $\Longleftrightarrow$
    

\[
    \lim_{z \to \infty} f(z) \ \text{există și este finită}.
    \]



    \item[(ii)] Infinitul este un pol $\Longleftrightarrow$
    

\[
    \lim_{z \to \infty} |f(z)| = \infty.
    \]



    \item[(iii)] Infinitul este o singularitate esențială $\Longleftrightarrow$
    $f$ nu este mărginită în nicio regiune de forma $\{ |z| > R \}$.
\end{enumerate}
\end{proposition}


\end{document}
